\section{Proofs for \texorpdfstring{Section~\ref{sec:obj}}{conflicting objectives, deliberation time and scale selection}}\label{app:obj}

Throughout, $E_i\in\R^{d_u\times d_i}$ is the inclusion of member $i$'s coordinates, so that $u_i=E_i^\top u$ and $\sum_iE_iK_{ii}^{-1}E_i^\top=K_{\mathrm d}^{-1}$, and $n\ge2$. In the setting of \cref{def:obj-quadratic}, $G^{\mathrm{act}}_i(u)=\frac12\Vert\mathsf z_i+\sum_j\mathsf M^{(i)}_ju_j\Vert^2_{\Lambda^\star_i}$, so as a function of $u_i$ at fixed $u_{-i}$,
\begin{equation}\label{eq:obj-conditional}
G^{\mathrm{act}}_i(u)=\tfrac12u_i^\top K_{ii}u_i+u_i^\top\Bigl(\sum_{j\ne i}K_{ij}u_j+\mathsf h_i\Bigr)+\mathsf c_i(u_{-i}),
\end{equation}
which gives \eqref{eq:obj-profile} and the Gaussian conditional $p_i(\cdot\mid u_{-i})\propto e^{-\gamma G^{\mathrm{act}}_i}$ of \cref{def:obj-quadratic}; $K_{ii}\succ0$ because $\mathsf M^{(i)}_i$ has full column rank. Write $T_i$ for the kernel that redraws $u_i$ from $p_i(\cdot\mid u_{-i})$ and $T=\frac1n\sum_iT_i$ for the heat-bath kernel. With $\xi\sim\Ncal(0,(\gamma K_{ii})^{-1})$ independent of $u$, the update of member $i$ is
\begin{equation}\label{eq:obj-update}
u^+=\mathsf A_iu-E_iK_{ii}^{-1}\mathsf h_i+E_i\xi,\qquad\mathsf A_i=I-E_iK_{ii}^{-1}E_i^\top K.
\end{equation}

\begin{proof}[Proof of \cref{thm:obj-rotation}]
(a) The mixed second difference of a quadratic function over displacements $\Delta_i$ and $\Delta_j$ of two different blocks of coordinates is $\Delta_i^\top H_{ij}\Delta_j$, with $H$ its Hessian. The Hessian block of $G^{\mathrm{act}}_i$ is $\mathsf M^{(i)\top}_i\Lambda^\star_i\mathsf M^{(i)}_j=K_{ij}$ and that of $G^{\mathrm{act}}_j$ is $\mathsf M^{(j)\top}_i\Lambda^\star_j\mathsf M^{(j)}_j=K_{ji}^\top$. Hence $\mathcal L_{ij}=\partial^2_{ij}G^{\mathrm{act}}_j-\partial^2_{ij}G^{\mathrm{act}}_i=\Delta_i^\top(K_{ji}^\top-K_{ij})\Delta_j$, which is \eqref{eq:obj-cycle} because $(K_{\mathrm a})_{ij}=(K_{ij}-K_{ji}^\top)/2$. The diagonal blocks of $K_{\mathrm a}$ vanish since each $K_{ii}$ is symmetric. If $K_{\mathrm a}=0$, the criterion \eqref{eq:agy-acts-criterion} holds on every finite restriction, which is then an exact potential game and reversible by \cref{thm:agy-joint-acts}(a) and \cref{thm:agy-agency}(b). If $(K_{\mathrm a})_{ij}\neq0$, choose $\Delta_i,\Delta_j$ with $\Delta_i^\top(K_{\mathrm a})_{ij}\Delta_j\neq0$; the restriction to $\mathcal U_i=\{0,\Delta_i\}$, $\mathcal U_j=\{0,\Delta_j\}$ and $\mathcal U_k=\{0\}$ otherwise has a four-cycle with ratio $e^{-\gamma\mathcal L_{ij}}\neq1$ and is not reversible.

(b) Sufficiency. Let $K=K^\top\succ0$, $\Phi(u)=\frac12u^\top Ku+\mathsf h^\top u$ and $\mu\propto e^{-\gamma\Phi}$, that is $\mu=\Ncal(u^{\mathrm N},(\gamma K)^{-1})$. By symmetry of $K$ and \eqref{eq:obj-conditional}, $\Phi-G^{\mathrm{act}}_i$ does not depend on $u_i$, so the conditional law of $u_i$ given $u_{-i}$ under $\mu$ is $p_i(\cdot\mid u_{-i})$. The measure $\mu(\mathrm du)T_i(u,\mathrm dv)$ then has, on $\{u_{-i}=v_{-i}\}$, the density $\mu_{-i}(u_{-i})p_i(u_i\mid u_{-i})p_i(v_i\mid u_{-i})$, which is symmetric in $(u_i,v_i)$. Each $T_i$, and so $T$, is reversible with respect to $\mu$. For uniqueness, run two copies with the same selected member and the same $\xi$. By \eqref{eq:obj-update} their difference obeys $e^+=\mathsf A_ie$, and $\mathsf A_i$ is the projector along $\operatorname{range}E_i$ that is orthogonal in the inner product of $K$, so $\Vert\mathsf A_ie\Vert_K^2=\Vert e\Vert_K^2-(E_i^\top Ke)^\top K_{ii}^{-1}(E_i^\top Ke)$ and
\begin{equation}\label{eq:obj-coupling}
\E\bigl[\Vert e^+\Vert_K^2\mid e\bigr]=\Vert e\Vert_K^2-\frac1n(Ke)^\top K_{\mathrm d}^{-1}(Ke)\le\Bigl(1-\frac\kappa n\Bigr)\Vert e\Vert_K^2,
\end{equation}
because $(Ke)^\top K_{\mathrm d}^{-1}Ke=w^\top K^{1/2}K_{\mathrm d}^{-1}K^{1/2}w$ with $w=K^{1/2}e$, and $K^{1/2}K_{\mathrm d}^{-1}K^{1/2}$ has the spectrum of $K_{\mathrm d}^{-1/2}KK_{\mathrm d}^{-1/2}$, with smallest eigenvalue $\kappa>0$. If $\nu$ is stationary, start one copy in $\nu$ and an independent one in $\mu$; then $e_t\to0$ in probability while the two marginal laws stay $\nu$ and $\mu$, so $\nu=\mu$.

Necessity. Let $\mu$ be a probability law with $\mathfrak M=\mu\otimes T$ invariant under the swap $(u,v)\mapsto(v,u)$, and put $\mathfrak M_i=\mu\otimes T_i$, $S_i=\{(u,v):u_{-i}=v_{-i},u_i\neq v_i\}$. The sets $S_i$ are swap invariant and pairwise disjoint; $\mathfrak M_i$ is carried by $S_i$, since $p_i$ has a density, and gives no mass to $S_j$, $j\neq i$. Hence $\frac1n\mathfrak M_i$ is the restriction of $\mathfrak M$ to $S_i$ and is swap invariant, so its two marginals agree: $\mu T_i=\mu$ for every $i$. It follows that under $\mu$ the conditional law of $u_i$ given $u_{-i}$ is $p_i(\cdot\mid u_{-i})$, and that $\mu=\mu T_1\cdots T_n$ has the density $m(v)=\int\mu(\mathrm du)\prod_ip_i(v_i\mid v_1,\dots,v_{i-1},u_{i+1},\dots,u_n)$, which is positive everywhere. For each $i$ there is therefore a measurable $c_i$ with
\begin{equation}\label{eq:obj-logdensity}
\ln m(u)=-\gamma G^{\mathrm{act}}_i(u)+c_i(u_{-i})\qquad\text{for almost every }u.
\end{equation}
Fix $i\neq j$ and displacements $\Delta_i,\Delta_j$, and let $\mathsf Df(u)=f(u+\Delta_i+\Delta_j)-f(u+\Delta_i)-f(u+\Delta_j)+f(u)$. Null sets are translation invariant, so for almost every $u$ the identity \eqref{eq:obj-logdensity} holds for $i$ and for $j$ at all four points. A function that does not depend on $u_i$, or does not depend on $u_j$, is annihilated by $\mathsf D$. Therefore $-\gamma\Delta_i^\top K_{ij}\Delta_j=\mathsf D\ln m(u)=-\gamma\Delta_i^\top K_{ji}^\top\Delta_j$ for all displacements, and $K_{ij}=K_{ji}^\top$: $K$ is symmetric. Then $\Phi-G^{\mathrm{act}}_i$ does not depend on $u_i$, and $h=\ln m+\gamma\Phi$ agrees almost everywhere with a function of $u_{-i}$, for every $i$. Let $g=\arctan h$ and let $\mathsf P_i$ integrate out $u_i$ against a standard Gaussian law; by Fubini's theorem $\mathsf P_i$ respects equality almost everywhere, and $\mathsf P_ig=g$ almost everywhere. Hence $g=\mathsf P_n\cdots\mathsf P_1g$ almost everywhere, and the right side is a constant. So $m$ is proportional to $e^{-\gamma\Phi}$. If $w^\top Kw\le0$ for some $w\neq0$, then $s\mapsto\Phi(u+sw)$ is concave for every $u$, $\int e^{-\gamma\Phi(u+sw)}\mathrm ds=\infty$, and $e^{-\gamma\Phi}$ is not integrable. Hence $K\succ0$ and $\mu=\Ncal(u^{\mathrm N},(\gamma K)^{-1})$.

(c) The solution is $u_t=e^{-Kt}u_0-\int_0^te^{-Kr}\mathsf h\mathrm dr+\sqrt{2/\gamma}\int_0^te^{-K(t-r)}\mathrm d\mathsf W_r$. If every eigenvalue of $K$ has positive real part, $e^{-Kt}\to0$ exponentially, $\Sigma=(2/\gamma)\int_0^\infty e^{-Kr}e^{-K^\top r}\mathrm dr$ converges, is positive definite, and satisfies $K\Sigma+\Sigma K^\top=-(2/\gamma)\int_0^\infty\frac{\mathrm d}{\mathrm dr}(e^{-Kr}e^{-K^\top r})\mathrm dr=(2/\gamma)I$. The solution of this equation is unique: a solution $X$ of $KX+XK^\top=0$ has $e^{-Kr}Xe^{-K^\top r}$ constant in $r$ and tending to zero. Given $u_0$, the law of $u_t$ is Gaussian with mean tending to $u^{\mathrm N}=-K^{-1}\mathsf h$ and covariance tending to $\Sigma$, so $\Ncal(u^{\mathrm N},\Sigma)$ is the unique stationary law. Conversely, let $\lambda$ be an eigenvalue with $\operatorname{Re}\lambda\le0$ and $w\neq0$ a left eigenvector, $w^\ast K=\lambda w^\ast$. Then $y_t=w^\ast u_t=a_t(u_0)+\mathsf N_t$, where $a_t$ is a function of $u_0$ and $\mathsf N_t=\sqrt{2/\gamma}\int_0^te^{-\lambda(t-r)}w^\ast\mathrm d\mathsf W_r$ is a centered complex Gaussian variable independent of $u_0$ with $\E\lvert\mathsf N_t\rvert^2=(2/\gamma)\lvert w\rvert^2\int_0^te^{-2r\operatorname{Re}\lambda}\mathrm dr\ge(2/\gamma)\lvert w\rvert^2t$. One of its real and imaginary parts has variance at least $\lvert w\rvert^2t/\gamma$, and a real variable $a+X$ with $X\sim\Ncal(0,\sigma^2)$ independent of $a$ has $\mathbb P(\lvert a+X\rvert\le R)\le2R/(\sigma\sqrt{2\pi})$. Under a stationary law the numbers $\mathbb P(\lvert\operatorname{Re}y_t\rvert\le R)$ and $\mathbb P(\lvert\operatorname{Im}y_t\rvert\le R)$ do not depend on $t$, and the smaller of them would be at most $2R\sqrt{\gamma}/(\lvert w\rvert\sqrt{2\pi t})$ for every $t$, hence zero for every $R$, which is impossible. So there is no stationary law.

Let now $\mathsf Q=K\Sigma-\gamma^{-1}I$. Then $\mathsf Q+\mathsf Q^\top=K\Sigma+\Sigma K^\top-(2/\gamma)I=0$ and $K=(\gamma^{-1}I+\mathsf Q)\Sigma^{-1}$. The Fokker--Planck equation is $\partial_tp=-\nabla\cdot(p\mathsf v)$ with current velocity $\mathsf v=-K(u-u^{\mathrm N})-\gamma^{-1}\nabla\ln p$; at the stationary density $\nabla\ln p=-\Sigma^{-1}(u-u^{\mathrm N})$ and $\mathsf v=-(K-\gamma^{-1}\Sigma^{-1})(u-u^{\mathrm N})=-\mathsf Q\Sigma^{-1}(u-u^{\mathrm N})$. Its mean square is $\tr(\mathsf Q\Sigma^{-1}\mathsf Q^\top)\ge0$, which vanishes only for $\mathsf Q=0$. Expanding $K\Sigma K^\top=(\gamma^{-1}I+\mathsf Q)\Sigma^{-1}(\gamma^{-1}I+\mathsf Q)^\top$ and using $\tr(\mathsf Q\Sigma^{-1})=0$ gives $\tr(K\Sigma K^\top)=\gamma^{-2}\tr\Sigma^{-1}+\tr(\mathsf Q\Sigma^{-1}\mathsf Q^\top)$ and $\tr K=\gamma^{-1}\tr\Sigma^{-1}$, which is \eqref{eq:obj-entropy}. The quantity $\gamma\E\lvert\mathsf v\rvert^2$ is what is usually called the entropy production rate of the diffusion; only its vanishing is used here. The stationary process is Gaussian with $\Cov(u_\tau,u_0)=C(\tau)=e^{-K\tau}\Sigma$ for $\tau\ge0$, and time reversal replaces $C(\tau)$ by $C(\tau)^\top$, so the process is reversible if and only if $C(\tau)$ is symmetric for all $\tau$. Differentiating at $\tau=0$, this forces $K\Sigma$ symmetric, that is $\mathsf Q=0$; conversely $\mathsf Q=0$ gives $K\Sigma=\Sigma K^\top$, hence $K^k\Sigma=\Sigma(K^\top)^k$ for all $k$ and $C(\tau)$ symmetric. Finally $\mathsf Q=0$ means $K=(\gamma\Sigma)^{-1}$, which is symmetric, and if $K$ is symmetric then $(\gamma K)^{-1}$ solves the Lyapunov equation, so $\Sigma=(\gamma K)^{-1}$ and $\mathsf Q=0$.
\end{proof}

\begin{lemma}[Moments of heat-bath deliberation]\label{lem:obj-moments}
In the setting of \cref{def:obj-quadratic} let $K$ be invertible, $\bar{\mathsf A}=\frac1n\sum_i\mathsf A_i=I-\frac1nK_{\mathrm d}^{-1}K$ and $\mathcal T(S)=\frac1n\sum_i\mathsf A_iS\mathsf A_i^\top$ on symmetric matrices. Under heat-bath deliberation the mean $\mathsf m_t$ and the second moment $S_t=\E(u_t-u^{\mathrm N})(u_t-u^{\mathrm N})^\top$ obey
\begin{equation}\label{eq:obj-moments}
\mathsf m_{t+1}-u^{\mathrm N}=\bar{\mathsf A}(\mathsf m_t-u^{\mathrm N}),\qquad S_{t+1}=\mathcal T(S_t)+\frac1{n\gamma}K_{\mathrm d}^{-1}.
\end{equation}
Every stationary law with a finite first moment has mean $u^{\mathrm N}$. A stationary law with finite second moments exists if and only if the spectral radius of $\mathcal T$ is below one, and its second moment is then the unique fixed point of \eqref{eq:obj-moments}.
\end{lemma}

\begin{proof}
Since $E_i^\top Ku^{\mathrm N}=-\mathsf h_i$, \eqref{eq:obj-update} reads $u^+-u^{\mathrm N}=\mathsf A_i(u-u^{\mathrm N})+E_i\xi$, and averaging over the selected member and over $\xi$ gives \eqref{eq:obj-moments}, with $\frac1n\sum_iE_i(\gamma K_{ii})^{-1}E_i^\top=(n\gamma)^{-1}K_{\mathrm d}^{-1}$. A stationary mean satisfies $K_{\mathrm d}^{-1}K(\mathsf m-u^{\mathrm N})=0$, so $\mathsf m=u^{\mathrm N}$. The map $\mathcal T$ and its adjoint $\mathcal T^\ast(X)=\frac1n\sum_i\mathsf A_i^\top X\mathsf A_i$ for the trace inner product preserve the cone of positive semidefinite matrices, which is closed, convex, pointed, self-dual and has nonempty interior. A linear map of a finite-dimensional space that preserves such a cone has its spectral radius $\mathsf r$ as an eigenvalue with an eigenvector in the cone (the finite-dimensional Krein--Rutman theorem). Let $X\succeq0$, $X\neq0$, with $\mathcal T^\ast X=\mathsf rX$. If a stationary second moment $S$ exists, pairing $S=\mathcal T(S)+(n\gamma)^{-1}K_{\mathrm d}^{-1}$ with $X$ gives $(1-\mathsf r)\tr(XS)=(n\gamma)^{-1}\tr(XK_{\mathrm d}^{-1})>0$, so $\mathsf r<1$. If $\mathsf r<1$, the fixed point $S_\infty=(n\gamma)^{-1}\sum_{k\ge0}\mathcal T^k(K_{\mathrm d}^{-1})$ exists and is unique, and from the point mass at $u^{\mathrm N}$ the second moments increase to $S_\infty$, so the laws $\nu_t$ are tight. The kernel $T$ maps bounded continuous functions to bounded continuous functions by \eqref{eq:obj-update}, so a weak limit point of the time averages of $\nu_t$ is stationary, and its second moment is finite by lower semicontinuity and equals $S_\infty$ by uniqueness.
\end{proof}

\begin{proof}[Proof of \cref{prop:obj-protocol}]
Langevin deliberation. A real $2\times2$ matrix has both eigenvalues in the open right half plane if and only if its trace and determinant are positive. Here $\tr K=K_{11}+K_{22}>0$ and $\det K=K_{11}K_{22}(1-\chi)$, so by \cref{thm:obj-rotation}(c) a stationary law exists if and only if $\chi<1$.

Heat-bath deliberation. Write the two updates as $u_1^+=a_1u_2+\beta_1+\xi_1$ and $u_2^+=a_2u_1+\beta_2+\xi_2$, with $a_1=-K_{12}/K_{11}$, $a_2=-K_{21}/K_{22}$, $\beta_i=-\mathsf h_i/K_{ii}$ and $\xi_i\sim\Ncal(0,(\gamma K_{ii})^{-1})$, so that $a_1a_2=\chi$. Let $\mu$ be a stationary law, with no moment assumption, and $\nu_1,\nu_2$ its marginals. The first marginal of $\mu T_2$ is $\nu_1$ and that of $\mu T_1$ is the law of $a_1U_2+\beta_1+\xi_1$ with $U_2\sim\nu_2$, so $\mu=\frac12\mu T_1+\frac12\mu T_2$ gives $\nu_1=\operatorname{law}(a_1U_2+\beta_1+\xi_1)$, and in the same way $\nu_2=\operatorname{law}(a_2U_1+\beta_2+\xi_2)$. For the characteristic functions $\varphi_i$ of $\nu_i$ this reads $\varphi_1(t)=e^{\mathrm it\beta_1}\varphi_2(a_1t)e^{-t^2/(2\gamma K_{11})}$ and $\varphi_2(t)=e^{\mathrm it\beta_2}\varphi_1(a_2t)e^{-t^2/(2\gamma K_{22})}$, hence
\begin{equation}\label{eq:obj-charfn}
\varphi_1(t)=e^{\mathrm it(\beta_1+a_1\beta_2)}\varphi_1(\chi t)e^{-\sigma_1^2t^2/2},\qquad\sigma_1^2=\frac1{\gamma K_{11}}+\frac{a_1^2}{\gamma K_{22}}>0.
\end{equation}
If $\lvert\chi\rvert\ge1$, iterating \eqref{eq:obj-charfn} $k$ times gives $\lvert\varphi_1(t)\rvert\le\exp(-\sigma_1^2t^2\sum_{j<k}\chi^{2j}/2)\le e^{-k\sigma_1^2t^2/2}$ for every $k$, so $\varphi_1(t)=0$ for $t\neq0$, which contradicts continuity at $\varphi_1(0)=1$. There is no stationary law. If $\lvert\chi\rvert<1$, the same iteration and $\varphi_1(\chi^kt)\to1$ give $\varphi_1(t)=\exp(\mathrm it\mathsf m_1-s_1t^2/2)$ with $\mathsf m_1=(\beta_1+a_1\beta_2)/(1-\chi)$ and $s_1=\sigma_1^2/(1-\chi^2)$: the marginals are $\nu_i=\Ncal(\mathsf m_i,s_i)$, uniquely. The law $\mu T_1$ depends on $\mu$ only through $\nu_2$ and is the Gaussian law $\mathsf G_1$ of $(a_1U_2+\beta_1+\xi_1,U_2)$, with mean $(\mathsf m_1,\mathsf m_2)$ and covariance $\bigl(\begin{smallmatrix}s_1&a_1s_2\\a_1s_2&s_2\end{smallmatrix}\bigr)$; likewise $\mu T_2=\mathsf G_2$ with the same mean and covariance $\bigl(\begin{smallmatrix}s_1&a_2s_1\\a_2s_1&s_2\end{smallmatrix}\bigr)$. So $\mu=\frac12\mathsf G_1+\frac12\mathsf G_2$ is the only candidate, and it is stationary because its marginals are $\nu_1$ and $\nu_2$, whence $\mu T_i=\mathsf G_i$. The mean satisfies $\mathsf m_1=a_1\mathsf m_2+\beta_1$ and $\mathsf m_2=a_2\mathsf m_1+\beta_2$, that is $K\mathsf m=-\mathsf h$, so $\mathsf m=u^{\mathrm N}$. An equal mixture of two Gaussian laws with one mean and different covariances is not Gaussian: along a direction in which the variances $v_1\neq v_2$ differ, the fourth central moment $3(v_1^2+v_2^2)/2$ exceeds the Gaussian value $3(v_1+v_2)^2/4$. The covariances agree if and only if $a_1s_2=a_2s_1$, and
\begin{equation*}
\gamma(1-\chi^2)(a_1s_2-a_2s_1)=\frac{a_1}{K_{22}}+\frac{a_1a_2^2}{K_{11}}-\frac{a_2}{K_{11}}-\frac{a_2a_1^2}{K_{22}}=(1-\chi)\frac{K_{21}-K_{12}}{K_{11}K_{22}},
\end{equation*}
which vanishes if and only if $K_{12}=K_{21}$.

(a) The matrix is realized by $\mathsf M_1=e_1$, $\mathsf M_2=e_2$ in $\R^2$, shared by both members, with $\Lambda^\star_1=\bigl(\begin{smallmatrix}1&c\\c&\lambda\end{smallmatrix}\bigr)$, $\Lambda^\star_2=\bigl(\begin{smallmatrix}\lambda&-c\\-c&1\end{smallmatrix}\bigr)$, $\lambda>c^2$, and $\mathsf z_i=0$: the members agree on the effects and disagree on the sign of the correlation they prefer between the two outcomes. Here $\chi=-c^2$. The eigenvalues of $K$ are $1\pm\mathrm ic$, $K+K^\top=2I$ shows that $\Sigma=\gamma^{-1}I$ solves the Lyapunov equation, $\mathsf Q=\gamma^{-1}\bigl(\begin{smallmatrix}0&c\\-c&0\end{smallmatrix}\bigr)$, and $\gamma\tr(\mathsf Q\Sigma^{-1}\mathsf Q^\top)=\gamma^2\tr(\mathsf Q\mathsf Q^\top)=2c^2$. For the heat bath, $\lvert\chi\rvert<1$ means $\lvert c\rvert<1$, $a_1=-c$, $a_2=c$, $\sigma_1^2=\sigma_2^2=(1+c^2)/\gamma$ and $s_1=s_2=(1+c^2)/(\gamma(1-c^4))=s$, which gives the two components. Their average covariance is $sI$: the acts are uncorrelated and dependent, with $\E u_1^2u_2^2=s^2(1+2c^2)$ against $s^2$ for independent acts. For the moments, $\mathsf A_1=\bigl(\begin{smallmatrix}0&-c\\0&1\end{smallmatrix}\bigr)$ and $\mathsf A_2=\bigl(\begin{smallmatrix}1&0\\c&0\end{smallmatrix}\bigr)$, so $\bar{\mathsf A}=\frac12\bigl(\begin{smallmatrix}1&-c\\c&1\end{smallmatrix}\bigr)$ has eigenvalues $(1\pm\mathrm ic)/2$ of modulus $\sqrt{1+c^2}/2$, below one exactly for $c^2<3$, and on $S=\bigl(\begin{smallmatrix}s_1&r\\r&s_2\end{smallmatrix}\bigr)$ the map $\mathcal T$ of \cref{lem:obj-moments} acts as $(s_1,s_2,r)\mapsto((s_1+c^2s_2)/2,(c^2s_1+s_2)/2,c(s_1-s_2)/2)$, with eigenvalues $(1+c^2)/2$, $(1-c^2)/2$ and $0$. In particular $\tr S_{t+1}=\frac{1+c^2}2\tr S_t+\gamma^{-1}$ by \eqref{eq:obj-moments}, so for $c^2\ge1$ the variance grows at least as $t/\gamma$ from every initial law with finite second moments, while for $c^2<3$ the mean converges to $u^{\mathrm N}=0$.

(b) With the columns $\mathsf M_1=(1,\alpha)^\top$ and $\mathsf M_2=(\alpha,1)^\top$, $K_{11}=\mathsf M_1^\top\Lambda^\star_1\mathsf M_1=\epsilon+\alpha^2=K_{22}$ and $K_{12}=\mathsf M_1^\top\Lambda^\star_1\mathsf M_2=\alpha(1+\epsilon)=\mathsf M_2^\top\Lambda^\star_2\mathsf M_1=K_{21}$. So $K$ is symmetric, the game is an exact potential game on every finite restriction by \cref{thm:obj-rotation}(a), and $\chi=\alpha^2(1+\epsilon)^2/(\epsilon+\alpha^2)^2$. For $\alpha>0$, $\chi>1$ if and only if $\alpha(1+\epsilon)>\epsilon+\alpha^2$, that is $(1-\alpha)(\alpha-\epsilon)>0$, that is $\epsilon<\alpha<1$. At $\alpha=0.5$, $\epsilon=0.1$, $K=\bigl(\begin{smallmatrix}0.35&0.55\\0.55&0.35\end{smallmatrix}\bigr)$ has eigenvalues $-0.2$ and $0.9$ and $\chi=0.3025/0.1225>1$; by the first part neither protocol has a stationary law. The potential $\frac12u^\top Ku+\mathsf h^\top u$ is unbounded below along $u_1=-u_2$: member $1$ weights the second outcome, which member $2$'s act moves most, and answers each move of member $2$ by a larger move of its own. With shared precision and effect models $K=\mathsf M^\top\Lambda^\star\mathsf M$ is a Gram matrix, so $0\le\chi\le1$ by the Cauchy--Schwarz inequality, with $\chi=1$ only for collinear effect directions; for $\mathsf M$ of full column rank $K\succ0$ and both protocols have the stationary law $\Ncal(u^{\mathrm N},(\gamma K)^{-1})$.
\end{proof}

\begin{proof}[Proof of \cref{prop:obj-cost}]
Here $K_{ij}=\mathsf M_i^\top\Lambda^\star\mathsf M_j$, so $K=\mathsf M^\top\Lambda^\star\mathsf M$ is symmetric and positive definite. By \cref{thm:obj-rotation}(b) the heat bath has the unique stationary law $\Ncal(u^{\mathrm N},(\gamma K)^{-1})$, and by \cref{thm:obj-rotation}(c) so has Langevin deliberation, since $\Sigma=(\gamma K)^{-1}$ solves the Lyapunov equation. With $G_\Sigma=\sum_iG^{\mathrm{act}}_i$ and $G^{\mathrm{act}}_i(u)=\frac12\Vert\mathsf z_i+\mathsf Mu\Vert^2_{\Lambda^\star}$, the identity $\sum_i\Vert\mathsf z_i+v\Vert^2=n\Vert\bar{\mathsf z}+v\Vert^2+\sum_i\Vert\mathsf z_i-\bar{\mathsf z}\Vert^2$ gives $G_\Sigma(u)=\frac n2\Vert\bar{\mathsf z}+\mathsf Mu\Vert^2_{\Lambda^\star}+\text{const}$, a quadratic function with Hessian $nK$ and minimizer $u^{\mathrm C}=-K^{-1}\mathsf M^\top\Lambda^\star\bar{\mathsf z}$. Hence $G_\Sigma(u)-\min G_\Sigma=\frac n2(u-u^{\mathrm C})^\top K(u-u^{\mathrm C})$ and
\begin{equation*}
\E_\mu G_\Sigma-\min G_\Sigma=\frac n2(u^{\mathrm N}-u^{\mathrm C})^\top K(u^{\mathrm N}-u^{\mathrm C})+\frac n2\tr\bigl(K(\gamma K)^{-1}\bigr).
\end{equation*}
The last term is $nd_u/(2\gamma)$. Block $i$ of $K(u^{\mathrm N}-u^{\mathrm C})$ is $-\mathsf h_i+\mathsf M_i^\top\Lambda^\star\bar{\mathsf z}=-\mathsf w_i$, so the first term is $\frac n2\mathsf w^\top K^{-1}\mathsf w$, which is \eqref{eq:obj-cost}. A single agent that samples $u\propto e^{-\gamma G_\Sigma}$ has the law $\Ncal(u^{\mathrm C},(\gamma nK)^{-1})$ and pays $\frac n2\tr(K(\gamma nK)^{-1})=d_u/(2\gamma)$.

For the group model, $\mathsf z_{\mathrm N}=-\mathsf Mu^{\mathrm N}$ satisfies $\mathsf M_j^\top\Lambda^\star\mathsf z_{\mathrm N}=-(Ku^{\mathrm N})_j=\mathsf h_j=\mathsf M_j^\top\Lambda^\star\mathsf z_j$. In \eqref{eq:agy-act-potential} with $\mathsf z_C=\mathsf z_{\mathrm N}$ the own-act terms $\langle\mathsf z_j-\mathsf z_{\mathrm N},\mathsf M_ju_j\rangle_{\Lambda^\star}$ therefore vanish, the potential is the group act part $\frac12\Vert\mathsf z_{\mathrm N}+\mathsf Mu\Vert^2_{\Lambda^\star}=\frac12(u-u^{\mathrm N})^\top K(u-u^{\mathrm N})$, and $\mu\propto e^{-\gamma G^{\mathrm{act}}_C}$ is the group posterior with flat habits. A displacement $\mathsf z$ gives the same group posterior if and only if $\mathsf M^\top\Lambda^\star\mathsf z=\mathsf h$, and $\mathsf z_{\mathrm N}$ is the only such displacement in the range of $\mathsf M$.

For two scalar acts write $\mathsf M_i=\mathsf m_ie_i$ with $\mathsf m_i>0$ and $\Vert e_i\Vert_{\Lambda^\star}=1$, and $\mathsf d=(\mathsf z_1-\mathsf z_2)/2=\mathsf z_1-\bar{\mathsf z}=-(\mathsf z_2-\bar{\mathsf z})$. Then $K=\mathsf D\mathsf R\mathsf D$ with $\mathsf D=\diag(\mathsf m_1,\mathsf m_2)$ and $\mathsf R=\bigl(\begin{smallmatrix}1&\cos\theta\\\cos\theta&1\end{smallmatrix}\bigr)$, and $\mathsf w=\mathsf D(\eta_1,-\eta_2)^\top$. Hence $\mathsf w^\top K^{-1}\mathsf w=(\eta_1,-\eta_2)\mathsf R^{-1}(\eta_1,-\eta_2)^\top=(\eta_1^2+2\cos\theta\eta_1\eta_2+\eta_2^2)/\sin^2\theta$, which is \eqref{eq:obj-cost-two} since $n/2=1$. As $\theta\to0$ the numerator tends to $4\eta_1^2$, so the bias diverges as $\sin^{-2}\theta$ unless the members agree along the common direction.
\end{proof}

\begin{proof}[Proof of \cref{thm:obj-time}]
Let $\mathsf R=K_{\mathrm d}^{-1/2}KK_{\mathrm d}^{-1/2}\succ0$, with identity diagonal blocks. Its eigenvalues $\lambda_k$ average to one, so $\kappa=\min_k\lambda_k\in(0,1]$, and $x^\top\mathsf Rx=\Vert\sum_i\mathsf R^{1/2}E_ix_i\Vert^2\le n\sum_i\Vert\mathsf R^{1/2}E_ix_i\Vert^2=n\Vert x\Vert^2$, so $\lambda_k\in[\kappa,n]$. The matrix $K^{1/2}K_{\mathrm d}^{-1}K^{1/2}$ has the same spectrum.

(a) If $W_{2,K}(\nu_0,\mu)=\infty$ there is nothing to prove. Otherwise take an optimal coupling $(u_0,u'_0)$ of $\nu_0$ and $\mu$ and run both copies with the same selected members and the same $\xi$. Then $u'_t\sim\mu$ for all $t$ and, by \eqref{eq:obj-coupling}, $W_{2,K}(\nu_t,\mu)^2\le\E\Vert u_t-u'_t\Vert_K^2\le(1-\kappa/n)^tW_{2,K}(\nu_0,\mu)^2$. By \eqref{eq:obj-moments} the mean error is multiplied at each event by $\bar{\mathsf A}=I-\frac1nK_{\mathrm d}^{-1}K$. Since $K^{1/2}\bar{\mathsf A}K^{-1/2}=I-\frac1nK^{1/2}K_{\mathrm d}^{-1}K^{1/2}$ is symmetric with eigenvalues $1-\lambda_k/n\in[0,1-\kappa/n]$, $\bar{\mathsf A}$ is self-adjoint in the inner product of $K$ with norm $1-\kappa/n$, attained on the eigenvector of the smallest $\lambda_k$.

(b) Put $x=(\gamma K)^{1/2}(u-u^{\mathrm N})$, so that $x\sim\Ncal(0,I)$ under $\mu$. By \eqref{eq:obj-update}, an update of member $i$ reads $x^+=\mathsf Q_ix+(\gamma K)^{1/2}E_i\xi$ with $\mathsf Q_i=K^{1/2}\mathsf A_iK^{-1/2}=I-K^{1/2}E_iK_{ii}^{-1}E_i^\top K^{1/2}$, the orthogonal projector onto the complement of $\operatorname{range}(K^{1/2}E_i)$, and the noise has covariance $I-\mathsf Q_i$. Hence $T_if(x)=\E f(\mathsf Q_ix+(I-\mathsf Q_i)\zeta)$ with $\zeta\sim\Ncal(0,I)$. For $\varepsilon_a(x)=\exp(a^\top x-\Vert a\Vert^2/2)$ this gives $T_i\varepsilon_a=\exp(a^\top\mathsf Q_ix-\Vert a\Vert^2/2+\Vert(I-\mathsf Q_i)a\Vert^2/2)=\varepsilon_{\mathsf Q_ia}$. Decompose $L^2(\Ncal(0,I))=\bigoplus_{k\ge0}\mathcal H_k$ into Wiener chaoses and let $\mathsf I_k$ be the map from symmetric $k$-tensors onto $\mathcal H_k$ with $\mathsf I_k(a^{\otimes k})=\Vert a\Vert^k\mathrm{He}_k(a^\top x/\Vert a\Vert)$, for which $\langle\mathsf I_k(a^{\otimes k}),\mathsf I_k(b^{\otimes k})\rangle=k!\langle a,b\rangle^k$ and $\varepsilon_{sa}=\sum_ks^k\mathsf I_k(a^{\otimes k})/k!$ in $L^2$. Comparing the coefficients of $s^k$ in $T_i\varepsilon_{sa}=\varepsilon_{s\mathsf Q_ia}$ gives $T_i\mathsf I_k(a^{\otimes k})=\mathsf I_k((\mathsf Q_ia)^{\otimes k})$: every $T_i$ preserves each chaos, and $T$ is unitarily equivalent on $\mathcal H_k$ to the restriction of $\frac1n\sum_i\mathsf Q_i^{\otimes k}$ to symmetric tensors. On $\mathcal H_0$, the constants, it is the identity. On $\mathcal H_1$ it is $\frac1n\sum_i\mathsf Q_i=I-\frac1nK^{1/2}K_{\mathrm d}^{-1}K^{1/2}$, with largest eigenvalue $1-\kappa/n$. For $k\ge2$, $\mathsf Q_i^{\otimes k}$ is a product of the commuting projectors $\mathsf Q_i\otimes I\otimes\dots\otimes I$, \dots, $I\otimes\dots\otimes I\otimes\mathsf Q_i$, hence a projector below each of them, and $0\preceq\frac1n\sum_i\mathsf Q_i^{\otimes k}\preceq(\frac1n\sum_i\mathsf Q_i)\otimes I\otimes\dots\otimes I\preceq(1-\kappa/n)I$. The spectrum of the self-adjoint operator $T$ on the complement of the constants therefore lies in $[0,1-\kappa/n]$ and contains $1-\kappa/n$.

(c) For two members $\mathsf R=\bigl(\begin{smallmatrix}I&\mathsf C\\\mathsf C^\top&I\end{smallmatrix}\bigr)$ with $\mathsf C=K_{11}^{-1/2}K_{12}K_{22}^{-1/2}$, whose eigenvalues are $1\pm\sigma_k(\mathsf C)$ and $1$, so $\kappa=1-\sigma_{\max}(\mathsf C)$. Under shared $\Lambda^\star$ and effect models, $\mathsf U_i=(\Lambda^\star)^{1/2}\mathsf M_iK_{ii}^{-1/2}$ has orthonormal columns spanning $\operatorname{range}((\Lambda^\star)^{1/2}\mathsf M_i)$ and $\mathsf C=\mathsf U_1^\top\mathsf U_2$, whose singular values are by definition the cosines of the principal angles between the two ranges.

(d) The convention is that at event $t$ the selected member responds to the loads $\mathsf h_t$, after which the loads move; $u^{\mathrm N}_t=-K^{-1}\mathsf h_t$ and $\mu_t=\Ncal(u^{\mathrm N}_t,(\gamma K)^{-1})$. With $\mathsf e_t=\E_{\nu_t}u-u^{\mathrm N}_t$, \eqref{eq:obj-moments} gives $\mathsf e_{t+1}=\bar{\mathsf A}\mathsf e_t-(u^{\mathrm N}_{t+1}-u^{\mathrm N}_t)$, so $\Vert\mathsf e_{t+1}\Vert_K\le(1-\kappa/n)\Vert\mathsf e_t\Vert_K+\delta$ and $\Vert\mathsf e_t\Vert_K\le(1-\kappa/n)^t\Vert\mathsf e_0\Vert_K+\delta\sum_{j<t}(1-\kappa/n)^j$, which is the first bound. For a constant drift $\mathsf v$ per event the error converges to $-(I-\bar{\mathsf A})^{-1}\mathsf v=-nK^{-1}K_{\mathrm d}\mathsf v$; if $\mathsf v$ is an eigenvector of $\bar{\mathsf A}$ for $1-\kappa/n$ with $\Vert\mathsf v\Vert_K=\delta$, this limit is $-(n/\kappa)\mathsf v$, of norm $n\delta/\kappa$, so the constant cannot be improved. For the laws, $\mu_t$ is stationary for the kernel $T^{(t)}$ with loads $\mathsf h_t$, the contraction of (a) applies to any pair of laws, and two Gaussian laws with equal covariances are at transport distance equal to the distance of their means. Hence
\begin{equation*}
W_{2,K}(\nu_{t+1},\mu_{t+1})\le W_{2,K}\bigl(\nu_tT^{(t)},\mu_tT^{(t)}\bigr)+W_{2,K}(\mu_t,\mu_{t+1})\le(1-\kappa/n)^{1/2}W_{2,K}(\nu_t,\mu_t)+\delta,
\end{equation*}
and the second bound follows from $1/(1-\sqrt{1-\kappa/n})\le2n/\kappa$.
\end{proof}

\begin{proposition}[Rotation does not slow Langevin deliberation]\label{prop:obj-langevin-rate}
Let every eigenvalue of $K$ have positive real part, let $\Sigma$ solve $K\Sigma+\Sigma K^\top=(2/\gamma)I$ and $K_0=(\gamma\Sigma)^{-1}$. Every eigenvalue $\lambda$ of $K$ satisfies $\lambda_{\min}(K_0)\le\operatorname{Re}\lambda\le\lambda_{\max}(K_0)$. Conversely, for every $\Sigma\succ0$ and antisymmetric $\mathsf Q$ the matrix $(\gamma^{-1}I+\mathsf Q)\Sigma^{-1}$ has all eigenvalues in the open right half plane.
\end{proposition}

\begin{proof}
Let $K=(\gamma^{-1}I+\mathsf Q)\Sigma^{-1}$ with $\mathsf Q$ real antisymmetric, as in \cref{thm:obj-rotation}(c) or as given, $Kx=\lambda x$ with $x\neq0$ complex, and $y=\Sigma^{-1}x$. Then $(\gamma^{-1}I+\mathsf Q)y=\lambda\Sigma y$, and multiplying by $y^\ast$ gives $\gamma^{-1}\Vert y\Vert^2+y^\ast\mathsf Qy=\lambda y^\ast\Sigma y$. The number $y^\ast\mathsf Qy$ is imaginary and $y^\ast\Sigma y>0$, so $\operatorname{Re}\lambda=\gamma^{-1}\Vert y\Vert^2/(y^\ast\Sigma y)$, which lies between $1/(\gamma\lambda_{\max}(\Sigma))$ and $1/(\gamma\lambda_{\min}(\Sigma))$.
\end{proof}

Among the interaction matrices with one stationary law the potential one therefore has the slowest asymptotic decay of means and lagged covariances; the heat bath behaves in the opposite way in \cref{prop:obj-protocol}(a), where rotation raises the modulus $\sqrt{1+c^2}/2$ of the mean map.

\begin{proof}[Proof of \cref{prop:obj-finite}]
Write $p_i(\cdot\mid\pi_{-i})\propto E_ie^{-\gamma G_i(\cdot,\pi_{-i})}$ for the conditionals and $d_{\mathrm H}$ for the Hamming distance. If $\pi_{-i}$ and $\pi'_{-i}$ differ only in member $j$'s policy, $p_i(\cdot\mid\pi'_{-i})\propto p_i(\cdot\mid\pi_{-i})e^g$ with $g=-\gamma(G_i(\cdot,\pi'_{-i})-G_i(\cdot,\pi_{-i}))$, of oscillation at most $\gamma\omega_{ij}$, and \cref{lem:agy-tilt} gives $\mathfrak J_{ij}\le\tanh(\gamma\omega_{ij}/4)$.

One event from adjacent profiles. Let $\pi$ and $\pi'$ differ only in member $j$'s policy, select the same member $i$ in both copies, and draw the two new policies from a maximal coupling of $p_i(\cdot\mid\pi_{-i})$ and $p_i(\cdot\mid\pi'_{-i})$. If $i=j$ the two conditionals coincide and the copies coalesce. If $i\neq j$ the copies still differ at $j$ and differ at $i$ with probability at most $\mathfrak J_{ij}$. Hence the expected Hamming distance after the event is at most $\frac1n\sum_{i\ne j}(1+\mathfrak J_{ij})\le1-\frac1n+\frac{\mathfrak D}n=\rho$.

One event from arbitrary laws. Two profiles at Hamming distance $k$ are joined by a path of $k$ adjacent steps, so by the triangle inequality $W_1(\delta_\pi T,\delta_{\pi'}T)\le\rho d_{\mathrm H}(\pi,\pi')$. For laws $\nu,\nu'$ with optimal coupling $\vartheta$, mixing optimal couplings of $\delta_\pi T$ and $\delta_{\pi'}T$ over $\vartheta$ gives a coupling of $\nu T$ and $\nu'T$, so $W_1(\nu T,\nu'T)\le\rho W_1(\nu,\nu')$. If $\mathfrak D<1$ the stationary law is unique and $W_1(\nu_t,\mathbb P^\infty)\le\rho^tW_1(\nu_0,\mathbb P^\infty)\le n\rho^t$, the diameter being $n$; and $\mathrm{TV}\le W_1$ because $d_{\mathrm H}(\pi,\pi')\ge\one[\pi\neq\pi']$. This is \eqref{eq:obj-finite-mixing}.

Two profiles of expected free energies. Let $T'$ be the kernel of $G'$ and $\mathbb P'^\infty$ a stationary law of it. From one profile $\pi$, select the same member $i$ and couple $p_i(\cdot\mid\pi_{-i})$ and $p'_i(\cdot\mid\pi_{-i})\propto p_i(\cdot\mid\pi_{-i})e^{-\gamma(G'_i-G_i)}$ maximally; by \cref{lem:agy-tilt} the copies differ, in one coordinate, with probability at most $\tanh(\gamma\omega_i/4)$. So $W_1(\delta_\pi T,\delta_\pi T')\le\frac1n\sum_i\tanh(\gamma\omega_i/4)=\tau$ for every $\pi$, and
\begin{equation*}
W_1\bigl(\mathbb P^\infty,\mathbb P'^\infty\bigr)\le W_1\bigl(\mathbb P^\infty T,\mathbb P'^\infty T\bigr)+W_1\bigl(\mathbb P'^\infty T,\mathbb P'^\infty T'\bigr)\le\rho W_1\bigl(\mathbb P^\infty,\mathbb P'^\infty\bigr)+\tau,
\end{equation*}
which gives \eqref{eq:obj-finite-shift} since $1-\rho=(1-\mathfrak D)/n$. Only the influence of one of the two chains enters.

Moving expected free energies. Let $T_t$ be the kernel of the profile $G_{\cdot,t}$ in force at event $t$, $\nu_{t+1}=\nu_tT_t$, $\mathfrak D_t\le\bar{\mathfrak D}<1$ its influence constant and $\mathbb P^\infty_t$ its stationary law, which is unique because $T_t$ contracts $W_1$, with $\omega_{i,t}=\max_{\pi_{-i}}\operatorname{osc}_{\pi_i}(G_{i,t+1}-G_{i,t})$. The bound \eqref{eq:obj-finite-shift}, applied to the profiles at events $t$ and $t+1$ with the influence constant $\mathfrak D_t$ of the first, gives $W_1(\mathbb P^\infty_t,\mathbb P^\infty_{t+1})\le\sum_i\tanh(\gamma\omega_{i,t}/4)/(1-\mathfrak D_t)\le n\bar\tau/(1-\bar{\mathfrak D})$, the kernel $T_t$ contracts $W_1$ by $1-(1-\mathfrak D_t)/n\le\bar\rho$, and $\mathbb P^\infty_tT_t=\mathbb P^\infty_t$, so
\begin{equation*}
W_1\bigl(\nu_{t+1},\mathbb P^\infty_{t+1}\bigr)\le W_1\bigl(\nu_tT_t,\mathbb P^\infty_tT_t\bigr)+W_1\bigl(\mathbb P^\infty_t,\mathbb P^\infty_{t+1}\bigr)\le\bar\rho W_1\bigl(\nu_t,\mathbb P^\infty_t\bigr)+\frac{n\bar\tau}{1-\bar{\mathfrak D}},
\end{equation*}
and summing the geometric series gives \eqref{eq:obj-finite-tracking}. If each member's expected free energy moves by at most $\dot\omega/n$ in oscillation per event, that is $\dot\omega$ per round of $n$ events, then $\tanh(\gamma\omega_{i,t}/4)\le\gamma\dot\omega/(4n)$, so $\bar\tau$ can be taken to be $\gamma\dot\omega/(4n)$, and $W_1/n$, the expected fraction of members whose policy differs from that of a coupled draw of the current stationary law, is at most $\gamma\dot\omega/(4(1-\bar{\mathfrak D})^2)$ once the transient has decayed. Since $n\rho^t\le ne^{-t(1-\mathfrak D)/n}$, the law is within $\epsilon$ of stationarity after $\ln(n/\epsilon)/(1-\mathfrak D)$ rounds.
\end{proof}

\begin{proof}[Proof of \cref{prop:obj-identification}]
(a) If Langevin deliberation with matrix $K$ has the stationary law $\Ncal(u^{\mathrm N},\Sigma)$, then $K=(\gamma^{-1}I+\mathsf Q)\Sigma^{-1}$ with $\mathsf Q$ antisymmetric by \cref{thm:obj-rotation}(c). Conversely, for antisymmetric $\mathsf Q$ the matrix $K=(\gamma^{-1}I+\mathsf Q)\Sigma^{-1}$ satisfies $K\Sigma+\Sigma K^\top=(2/\gamma)I$ and has its spectrum in the open right half plane by \cref{prop:obj-langevin-rate}, so with loads $\mathsf h=-Ku^{\mathrm N}$ its stationary law is $\Ncal(u^{\mathrm N},\Sigma)$. The matrix $K$ is symmetric if and only if $\mathsf Q\Sigma^{-1}+\Sigma^{-1}\mathsf Q=0$, that is $\Sigma\mathsf Q+\mathsf Q\Sigma=0$, which forces $\mathsf Q=0$ because the spectra of $\Sigma$ and $-\Sigma$ are disjoint. For realizability, a matrix $K$ whose diagonal blocks are symmetric and positive definite is the interaction matrix of a quadratic deliberation: with outcome dimension at least $\max_id_i$ and isometries $\mathsf V_i$ of $\R^{d_i}$ into the outcome space, take $\Lambda^\star_i=I$, $\mathsf M^{(i)}_i=\mathsf V_iK_{ii}^{1/2}$ and $\mathsf M^{(i)}_j=\mathsf V_iK_{ii}^{-1/2}K_{ij}$, and $\mathsf z_i=\mathsf V_iK_{ii}^{-1/2}\mathsf h_i$. The diagonal blocks of $(\gamma^{-1}I+\mathsf Q)\Sigma^{-1}$ are $\gamma^{-1}(\Sigma^{-1})_{ii}+(\mathsf Q\Sigma^{-1})_{ii}$; for scalar acts they are positive for all small $\mathsf Q$, and for vector acts the realizable small $\mathsf Q$ are those for which every $(\mathsf Q\Sigma^{-1})_{ii}$ is symmetric.

(b) The change of variable $u\mapsto u-u^{\mathrm N}$ maps the deliberation with loads $\mathsf h$ to the one with loads zero under either protocol, by \eqref{eq:obj-update}, so the existence of a stationary law does not depend on $\mathsf h$. A Langevin stationary law is $\Ncal(u^{\mathrm N},\Sigma)$ by \cref{thm:obj-rotation}(c), and a heat-bath stationary law with a finite first moment has mean $u^{\mathrm N}$ by \cref{lem:obj-moments}; in both cases the mean is $u^{\mathrm N}=-K^{-1}\mathsf h$. Hence $-\partial\E u/\partial\mathsf h=K^{-1}$, which is symmetric if and only if $K$ is.

(c) For $\tau\ge0$, $\Cov(u_\tau,u_0)=e^{-K\tau}\Sigma$, whose derivative at $\tau=0$ is $-K\Sigma$. The symmetry statement is in the proof of \cref{thm:obj-rotation}(c), and $K^{-1}=\Sigma(\gamma^{-1}I+\mathsf Q)^{-1}=\gamma\Sigma(I+\gamma\mathsf Q)^{-1}$, where $I+\gamma\mathsf Q$ is invertible because $\mathsf Q$ has imaginary spectrum.
\end{proof}

The proofs for \cref{sec:obj-scale} rest on one lemma about Gaussian pairs, which is used at a stationary horizon in \cref{thm:obj-selection} and between two arbitrary times in \cref{prop:obj-clocks}. Recall $\iota(s)=-\frac12\ln(1-e^{-2s})$ on $(0,\infty)$, with $\iota(\infty)=0$. Since $e^{-2s}+e^{-2\iota(s)}=1$ is symmetric in $s$ and $\iota(s)$, $\iota$ is a decreasing involution of $(0,\infty)$, and $\iota(\theta)\ge c$ if and only if $\theta\le\iota(c)=\varsigma(c)$.

\begin{lemma}[Information retained by a frame]\label{lem:obj-compression}
Let $(\mathsf x,\mathsf x')$ be jointly Gaussian with $\mathsf x\in\R^N$, $S=\Cov(\mathsf x)\succ0$ and $\Cov(\mathsf x\mid\mathsf x')=S^{1/2}(I-\mathsf H)S^{1/2}$, where $\mathsf H$ is symmetric with orthonormal eigenvectors $\mathsf v_1,\dots,\mathsf v_N$ and eigenvalues $e^{-2\theta_1}\ge\dots\ge e^{-2\theta_N}$, $0<\theta_1\le\dots\le\theta_N\le\infty$. Put $\theta_0=0$ and $\theta_{N+1}=\infty$.
\begin{enumerate}
\item[(a)] $I(\mathsf x;\mathsf x')=\sum_k\iota(\theta_k)$. For every $\Phi\in\R^{N\times M}$ with orthonormal columns, $I(\Phi^\top\mathsf x;\mathsf x')\le\sum_{k\le M}\iota(\theta_k)$, with equality if $S^{1/2}\operatorname{range}\Phi$ is spanned by $\mathsf v_1,\dots,\mathsf v_M$ and, if $\theta_M<\theta_{M+1}$, only then.
\item[(b)] For $c>0$ let $\mathsf j(M)=\sum_{k>M}\iota(\theta_k)+cM$, which by (a) is the minimum of $I(\mathsf x;\mathsf x')-I(\Phi^\top\mathsf x;\mathsf x')+c\rank\Phi$ over the frames of rank $M$. A rank $M\in\{0,\dots,N\}$ minimizes $\mathsf j$ if and only if $\theta_M\le\varsigma(c)\le\theta_{M+1}$, it is the only minimizing rank if and only if both inequalities are strict, and $M^\ast=\#\{k:\theta_k<\varsigma(c)\}$ is always a minimizing rank.
\end{enumerate}
\end{lemma}

\begin{proof}
(a) Let $\Psi=S^{1/2}\Phi$, which has rank $M$, and write $\Psi=W\mathsf R_0$ with $W\in\R^{N\times M}$ having orthonormal columns and $\mathsf R_0$ invertible, so that $\operatorname{range}W=S^{1/2}\operatorname{range}\Phi$. Then $\Cov(\Phi^\top\mathsf x)=\Psi^\top\Psi$ and $\Cov(\Phi^\top\mathsf x\mid\mathsf x')=\Psi^\top(I-\mathsf H)\Psi$, and the Gaussian formula $I=\frac12\ln\det\Cov(\Phi^\top\mathsf x)-\frac12\ln\det\Cov(\Phi^\top\mathsf x\mid\mathsf x')$ gives
\[
I(\Phi^\top\mathsf x;\mathsf x')=-\tfrac12\ln\det\bigl(W^\top(I-\mathsf H)W\bigr).
\]
The eigenvalues of $I-\mathsf H$ are $\mu_k=1-e^{-2\theta_k}$, with $0<\mu_1\le\dots\le\mu_N\le1$ and $-\frac12\ln\mu_k=\iota(\theta_k)$; the case $\Phi=I$ is the first claim. Let $\beta_1\le\dots\le\beta_M$ be the eigenvalues of $W^\top(I-\mathsf H)W$. By the Courant--Fischer principle, $\beta_k$ is the minimum over the $k$-dimensional subspaces of $\operatorname{range}W$ of the largest value of the quadratic form of $I-\mathsf H$ on unit vectors of the subspace, and $\mu_k$ is the same minimum over all $k$-dimensional subspaces of $\R^N$, so $\beta_k\ge\mu_k$ and $\det(W^\top(I-\mathsf H)W)\ge\prod_{k\le M}\mu_k$, which is the inequality. If $\operatorname{range}W$ is spanned by $\mathsf v_1,\dots,\mathsf v_M$, then $W^\top(I-\mathsf H)W$ is orthogonally similar to $\diag(\mu_1,\dots,\mu_M)$ and equality holds. Conversely, equality forces $\beta_k=\mu_k$ for every $k\le M$, hence $\tr(W^\top(I-\mathsf H)W)=\sum_{k\le M}\mu_k$. With $p_j=\Vert W^\top\mathsf v_j\Vert^2\in[0,1]$ and $\sum_jp_j=M$,
\[
\tr\bigl(W^\top(I-\mathsf H)W\bigr)-\sum_{k\le M}\mu_k=\sum_{j>M}\mu_jp_j-\sum_{k\le M}\mu_k(1-p_k)\ge(\mu_{M+1}-\mu_M)\sum_{j>M}p_j,
\]
because $\sum_{k\le M}(1-p_k)=\sum_{j>M}p_j$. If $\theta_M<\theta_{M+1}$ the factor $\mu_{M+1}-\mu_M$ is positive, so $p_j=0$ for $j>M$: the range of $W$ is orthogonal to $\mathsf v_{M+1},\dots,\mathsf v_N$ and, having dimension $M$, equals the span of $\mathsf v_1,\dots,\mathsf v_M$.

(b) $\mathsf j(M)-\mathsf j(M-1)=c-\iota(\theta_M)$ is nondecreasing in $M$ because $\iota$ decreases, so $\mathsf j$ is convex on $\{0,\dots,N\}$ and $M$ minimizes it if and only if $\mathsf j(M)-\mathsf j(M-1)\le0\le\mathsf j(M+1)-\mathsf j(M)$, each inequality being required when the difference is defined. This reads $\iota(\theta_M)\ge c\ge\iota(\theta_{M+1})$, that is $\theta_M\le\varsigma(c)\le\theta_{M+1}$. If one of the two is an equality, the neighboring rank has the same value; if both are strict, convexity makes $M$ the only minimizer. By its definition $M^\ast$ satisfies $\theta_{M^\ast}<\varsigma(c)\le\theta_{M^\ast+1}$.
\end{proof}

\begin{proof}[Proof of \cref{thm:obj-selection}]
In the stationary law $\Cov(x_0)=\Cov(x_\tau)=\Pi^{-1}$ and $x_\tau=e^{-\Pi\tau}x_0+\zeta$ with $\zeta$ Gaussian and independent of $x_0$, so $\Cov(x_0,x_\tau)=\Pi^{-1}e^{-\Pi\tau}$ and
\[
\Cov(x_0\mid x_\tau)=\Pi^{-1}-\Pi^{-1}e^{-\Pi\tau}\Pi e^{-\Pi\tau}\Pi^{-1}=\Pi^{-1/2}\bigl(I-e^{-2\Pi\tau}\bigr)\Pi^{-1/2}.
\]
\Cref{lem:obj-compression} applies with $S=\Pi^{-1}$, $\mathsf H=e^{-2\Pi\tau}$, $\theta_k=\varpi_k\tau$ and $\mathsf v_k$ the eigenvectors of $\Pi$.

(a) This is part (a) of the lemma, since $S^{1/2}\operatorname{range}\Phi=\Pi^{-1/2}\operatorname{range}\Phi$ equals $\operatorname{range}\mathsf V_{\mathrm s}$ if and only if $\operatorname{range}\Phi$ does, the slow eigenspace being invariant under $\Pi^{1/2}$.

(b) By (a), the minimum of $J_{\tau,c}$ over the frames of rank $M$ is $\mathsf j(M)$, attained on the frames with range $\operatorname{range}\mathsf V_{\mathrm s}$, so the minimum over all frames is $\min_M\mathsf j(M)$. Part (b) of the lemma gives \eqref{eq:obj-rank} and its uniqueness when no $\varpi_k\tau$ equals $\varsigma(c)$. If $\varsigma(c)/\varpi_{M+1}<\tau<\varsigma(c)/\varpi_M$, then $\varpi_M<\varpi_{M+1}$, $M$ is the only minimizing rank, and by (a) a frame of rank $M$ attains $\mathsf j(M)$ only if its range is the slow eigenspace. Conversely, if the minimizers are exactly the frames of rank $M$ on the slow eigenspace, then $M$ is the only minimizing rank, and both inequalities $\varpi_M\tau<\varsigma(c)<\varpi_{M+1}\tau$ are strict. The interval \eqref{eq:obj-windows} is nonempty exactly when $\varpi_M<\varpi_{M+1}$, and for $1\le M<N$ the logarithms of its endpoints differ by $\ln(\varpi_{M+1}/\varpi_M)$, in which $c$ does not appear.

(c) If $\operatorname{range}\Phi$ is invariant, $\Pi$ is block diagonal in the decomposition $\operatorname{range}\Phi\oplus(\operatorname{range}\Phi)^\perp$, with blocks $\Pi_\Phi$ and $\Pi_{\perp\perp}$, so its spectrum is the union of theirs. With $\varkappa_\perp>\varkappa_\Phi$ the $M$ smallest eigenvalues are those of $\Pi_\Phi$: $\operatorname{range}\Phi=\operatorname{range}\mathsf V_{\mathrm s}$, $\varpi_M=\varkappa_\Phi$ and $\varpi_{M+1}=\varkappa_\perp$. By (a), $J_{\tau,c}(\Phi)=\mathsf j(M)$, so $\Phi$ minimizes over all frames if and only if $M$ is a minimizing rank, that is $\varkappa_\Phi\tau\le\varsigma(c)\le\varkappa_\perp\tau$. For an equitable partition with uniform leak the span of the block indicators is invariant (\cref{lem:ren-galerkin}(a)), and, with $\xin^{\min}$ the minimum over the blocks with at least two members (one exists because $\rank\Phi<N$), \cref{thm:ren-memory}(d) gives $\varkappa_\perp\ge\xin^{\min}$ and $0<\varkappa_\Phi\le2\xout^{\max}-\varepsilon^\Pi$. Under $\xin^{\min}>2\xout^{\max}-\varepsilon^\Pi$ the interval $[\varsigma(c)/\xin^{\min},\varsigma(c)/(2\xout^{\max}-\varepsilon^\Pi)]$ is therefore contained in $[\varsigma(c)/\varkappa_\perp,\varsigma(c)/\varkappa_\Phi]$.
\end{proof}

For $\Pi=\varepsilon^\Pi I+L$ the eigenvalues are $\varpi_k=\varepsilon^\Pi+\lambda_k(L)$, so the condition $\varpi_k<\varsigma(c)/\tau$ of \eqref{eq:obj-rank} is $\lambda_k(L)<\varsigma(c)/\tau-\varepsilon^\Pi$, the strict mode rule of \cref{sec:ren-network} at $1/\tau_{\mathrm L}=\varsigma(c)/\tau-\varepsilon^\Pi$. A mode with $\varpi_k\tau=\varsigma(c)$ can be kept or discarded at no change of the criterion (\cref{lem:obj-compression}(b)), so the convention at equality does not affect optimality. Since $\varsigma(\iota(1))=1$, the price $c=\iota(1)=-\frac12\ln(1-e^{-2})\approx0.0727$ makes $\tau_{\mathrm L}=\tau$ at zero leak, and for $\tau\ge\varsigma(c)/\varepsilon^\Pi$ the right-hand side is not positive and nothing is retained. The reading of the diffusion time as a prediction horizon and the optimality statement are ours and concern the diffusion of \cref{thm:ren-memory}; we do not attribute them to \citet{villegas2023laplacian}.

\begin{remark}[The closure defect does not select a scale]\label{rem:obj-degenerate}
Let a criterion have the form $\mathsf f(\mathcal D)+c\rank\Phi$ with $\mathsf f\ge0$ and $\mathsf f(0)=0$, or the same with the memory kernel $\Gamma$ of \cref{thm:ren-memory}(b) in place of $\mathcal D$. Both $\mathcal D$ and $\Gamma$ vanish on every frame with invariant range (\cref{lem:ren-galerkin}(a); $\Pi_{\perp\Phi}=0$). The empty frame is one, so it attains the least possible value $0$ and the criterion selects nothing at every price. Restricted to a fixed rank $M\ge1$ the criterion is minimized by every frame spanned by $M$ eigenvectors of $\Pi$, the $M$ fastest included, although by \cref{thm:obj-selection}(a) these retain the least predictive information among invariant frames. In the nested medium every level has $\mathcal D=0$ (\cref{thm:ren-nested}(a)), so the defect is silent about which level to use. The closure defect measures whether a description is self-contained, which is a different property from being worth keeping, and this is why \eqref{eq:obj-criterion} uses predictive information. The statement is exact; the reading of the two properties is ours.
\end{remark}

\begin{proof}[Proof of \cref{thm:obj-nested}]
Let $\mathsf P_l$, $0\le l\le n_{\mathrm H}$, be the orthogonal projection of $\R^N$ onto the vectors that are constant on the level-$l$ blocks, $(\mathsf P_l)_{ij}=b^{-l}$ if $i$ and $j$ lie in one level-$l$ block and $0$ otherwise, so that $\mathsf P_0=I$, $\rank\mathsf P_l=b^{n_{\mathrm H}-l}$ and $\mathsf P_l\mathsf P_{l'}=\mathsf P_{\max(l,l')}$. The indicator of the pairs whose smallest common block has level $l\ge1$ is $b^l\mathsf P_l-b^{l-1}\mathsf P_{l-1}$, and every agent has $b^l-b^{l-1}$ such partners, so
\[
L=\sum_{l=1}^{n_{\mathrm H}}a^{\mathrm H}_l\bigl[(b^l-b^{l-1})I-b^l\mathsf P_l+b^{l-1}\mathsf P_{l-1}\bigr].
\]
Put $\mathsf E_m=\mathsf P_{m-1}-\mathsf P_m$ for $1\le m\le n_{\mathrm H}$ and $\mathsf E_{n_{\mathrm H}+1}=\mathsf P_{n_{\mathrm H}}$. These are orthogonal projections with $\mathsf E_m\mathsf E_{m'}=0$ for $m\ne m'$, they sum to $I$, $\rank\mathsf E_m=(b-1)b^{n_{\mathrm H}-m}$ for $m\le n_{\mathrm H}$, and $\mathsf P_l\mathsf E_m=\mathsf E_m$ if $l<m$ and $0$ otherwise. On $\operatorname{range}\mathsf E_m$ the bracket therefore equals $0$ for $l<m$, $b^m$ for $l=m$ and $b^l-b^{l-1}$ for $l>m$, so
\begin{equation}\label{eq:obj-nested-spectrum}
L=\sum_{m=1}^{n_{\mathrm H}}\mathsf g_m\mathsf E_m,\qquad\mathsf g_m=b^ma^{\mathrm H}_m+\sum_{l>m}(b^l-b^{l-1})a^{\mathrm H}_l,\qquad\mathsf g_m-\mathsf g_{m+1}=b^m\bigl(a^{\mathrm H}_m-a^{\mathrm H}_{m+1}\bigr),
\end{equation}
with $a^{\mathrm H}_{n_{\mathrm H}+1}=0$ and $\mathsf g_{n_{\mathrm H}+1}=0$. These identities hold for arbitrary level couplings. For $a^{\mathrm H}_l=A_{\mathrm h}\mathfrak q^{l-1}/((b-1)b^{l-1})$ one has $(b^l-b^{l-1})a^{\mathrm H}_l=A_{\mathrm h}\mathfrak q^{l-1}$ and $b^ma^{\mathrm H}_m=A_{\mathrm h}\mathfrak q^{m-1}+A_{\mathrm h}\mathfrak q^{m-1}/(b-1)$, so $\mathsf g_m$ is the quantity of \cref{prop:abs-hierarchy}, and $\mathsf g_m-\mathsf g_{m+1}=A_{\mathrm h}\mathfrak q^{m-1}(b-\mathfrak q)/(b-1)$ for $m<n_{\mathrm H}$, positive exactly when $\mathfrak q<b$.

(a) By \eqref{eq:obj-nested-spectrum}, $\Pi=\varepsilon^\Pi\mathsf E_{n_{\mathrm H}+1}+\sum_m(\varepsilon^\Pi+\mathsf g_m)\mathsf E_m$ with $\mathsf g_1>\dots>\mathsf g_{n_{\mathrm H}}>0$, which is the spectrum claimed. The number of eigenvalues not exceeding $\varepsilon^\Pi+\mathsf g_{l+1}$ is $1+\sum_{m>l}(b-1)b^{n_{\mathrm H}-m}=b^{n_{\mathrm H}-l}$, and the corresponding eigenspace is $\bigoplus_{m>l}\operatorname{range}\mathsf E_m=\operatorname{range}\mathsf P_l$, the span of the level-$l$ block indicators. Hence for $M=b^{n_{\mathrm H}-l}$, $0\le l\le n_{\mathrm H}$, one has $\varpi_M=\varepsilon^\Pi+\mathsf g_{l+1}=1/\tau_{l+1}$ and $\varpi_{M+1}=\varepsilon^\Pi+\mathsf g_l=1/\tau_l$ (with $\varpi_{N+1}=\infty$ and $\tau_0=0$ for $l=0$), and \cref{thm:obj-selection}(b) gives the windows $\varsigma(c)\tau_l<\tau<\varsigma(c)\tau_{l+1}$. For $M=0$ the window is $\tau>\varsigma(c)/\varpi_1=\varsigma(c)/\varepsilon^\Pi$. Every other rank has $\varpi_M=\varpi_{M+1}$ and by \cref{lem:obj-compression}(b) minimizes only at the single horizon $\tau=\varsigma(c)/\varpi_M$.

(b) $\sum_{l'=l}^{n_{\mathrm H}}\mathfrak q^{l'-1}=(\mathfrak q^{l-1}-\mathfrak q^{n_{\mathrm H}})/(1-\mathfrak q)$, so $\mathsf g_l=\mathsf g_1^\infty\mathfrak q^{l-1}-A_{\mathrm h}\mathfrak q^{n_{\mathrm H}}/(1-\mathfrak q)$ with $\mathsf g_1^\infty=A_{\mathrm h}(1/(1-\mathfrak q)+1/(b-1))$, which converges to $\mathsf g_1^\infty\mathfrak q^{l-1}$ as $n_{\mathrm H}\to\infty$. Since $\lvert C_l\rvert^{z_{\mathrm H}}=b^{l\ln(1/\mathfrak q)/\ln b}=\mathfrak q^{-l}$, the limit of $\tau_l$ is $1/(\varepsilon^\Pi+\mathsf g_1^\infty/(\mathfrak q\lvert C_l\rvert^{z_{\mathrm H}}))$, which is \eqref{eq:obj-scale-time}. If level-$l$ blocks have linear size proportional to $b^{l/D_e}$, then $\lvert C_l\rvert^{z_{\mathrm H}}$ is that size to the power $D_ez_{\mathrm H}=y_{\mathrm{RG}}$ (\cref{thm:ren-nested}(b)). The second term of the denominator of \eqref{eq:obj-scale-time} relative to the first is $\varepsilon^\Pi\mathfrak q^{1-l}/\mathsf g_1^\infty=\mathfrak x_lA_{\mathrm h}\mathfrak q/\mathsf g_1^\infty$, so the saturation at $1/\varepsilon^\Pi$ sets in where $\mathfrak x_l$ is of order one.

(c) By \eqref{eq:obj-nested-spectrum}, $e^{-\tau_{\mathrm L}L}=\sum_{m=1}^{n_{\mathrm H}}e^{-\tau_{\mathrm L}\mathsf g_m}(\mathsf P_{m-1}-\mathsf P_m)+\mathsf P_{n_{\mathrm H}}$. For a pair whose smallest common block has level $m_0$ ($m_0=0$ on the diagonal), $(\mathsf P_l)_{ij}=b^{-l}$ for $l\ge m_0$ and $0$ for $l<m_0$, so the entry is
\[
\mathsf k(m_0)=b^{-n_{\mathrm H}}+\sum_{m>m_0}e^{-\tau_{\mathrm L}\mathsf g_m}\bigl(b^{-(m-1)}-b^{-m}\bigr)-\one\{m_0\ge1\}e^{-\tau_{\mathrm L}\mathsf g_{m_0}}b^{-m_0},
\]
a function of the level alone. For $0\le m<n_{\mathrm H}$, $\mathsf k(m)-\mathsf k(m+1)=b^{-m}(e^{-\tau_{\mathrm L}\mathsf g_{m+1}}-e^{-\tau_{\mathrm L}\mathsf g_m})$ with $e^{-\tau_{\mathrm L}\mathsf g_0}=0$, which is positive because $\mathsf g_{m+1}<\mathsf g_m$, and $\mathsf k(n_{\mathrm H})=b^{-n_{\mathrm H}}(1-e^{-\tau_{\mathrm L}\mathsf g_{n_{\mathrm H}}})>0$. The density matrix is this kernel divided by its trace. The mode rule of \cref{sec:ren-network} keeps the eigenvalues strictly below $1/\tau_{\mathrm L}$ and discards a mode at equality; under the opposite convention at equality each window would lose its left endpoint and gain its right one. For $1/\mathsf g_l\le\tau_{\mathrm L}<1/\mathsf g_{l+1}$ the kept eigenvalues are $0$ and the $\mathsf g_m$ with $m>l$, in number $b^{n_{\mathrm H}-l}$. Pairs are joined in the order of increasing level, in an arbitrary order within a level, until $b^{n_{\mathrm H}-l}$ clusters remain. When all pairs of level at most $m$ have been joined, the clusters are the level-$m$ blocks, $b^{n_{\mathrm H}-m}$ in number. While pairs of level $m$ are being joined, each join merges two clusters inside one level-$m$ block or does nothing, so the clusters refine the level-$m$ blocks, and their number is at least $b^{n_{\mathrm H}-m}$ with equality only when they are the level-$m$ blocks. Since the number of clusters at the end of level $l-1$ is $b^{n_{\mathrm H}-l+1}>b^{n_{\mathrm H}-l}$, it first equals $b^{n_{\mathrm H}-l}$ during level $l$, at which moment the clusters are the level-$l$ blocks, whatever the order within the level. For $l=0$ no pair is joined.

Let $C$ be a level-$l$ block, $1\le l<n_{\mathrm H}$. Its induced medium is the nested medium of depth $l$ with the same couplings, so by \eqref{eq:obj-nested-spectrum} at depth $l$, $\lambda_2(L_C)=b^la^{\mathrm H}_l=bA_{\mathrm h}\mathfrak q^{l-1}/(b-1)$, and $\cut(C)/\lvert C\rvert=\sum_{l'>l}(b^{l'}-b^{l'-1})a^{\mathrm H}_{l'}=\sum_{l'>l}A_{\mathrm h}\mathfrak q^{l'-1}$. Hence
\[
\mathsf g_l=\lambda_2(L_C)+\frac{\cut(C)}{\lvert C\rvert},\qquad\mathsf g_{l+1}=\frac{\cut(C)}{\lvert C\rvert}+\frac{\mathfrak q}{b}\lambda_2(L_C).
\]
The block $C$ is a macronode exactly for $\tau_{\mathrm L}\in[1/\mathsf g_l,1/\mathsf g_{l+1})$. Since $1/\mathsf g_l<1/\lambda_2(L_C)$ and $1/\mathsf g_{l+1}<\lvert C\rvert/\cut(C)$, this window meets the interval $(1/\lambda_2(L_C),\lvert C\rvert/\cut(C))$ of \cref{conj:ren-lrg} if and only if $\mathsf g_{l+1}<\lambda_2(L_C)$, that is $\cut(C)/\lvert C\rvert<(1-\mathfrak q/b)\lambda_2(L_C)$, that is $r_\star(C)>b/(b-\mathfrak q)$.
\end{proof}

\begin{proof}[Proof of \cref{prop:obj-lrg-counter}]
Let $b\ge3$. The induced graph on $C$ is complete on $b-1$ vertices, so $\lambda_2(L_C)=b-1$; each vertex of $C$ has one edge to the remaining vertex, so $\cut(C)=\lvert C\rvert=b-1$ and $r_\star(C)=b-1$. The Laplacian of the complete graph has eigenvalues $0$ and $b$, the latter of multiplicity $b-1$. For $\tau_{\mathrm L}<1/b$ all $b$ modes are kept and the macronodes are the singletons; for $\tau_{\mathrm L}\ge1/b$ one mode is kept and the single macronode is the whole graph (the strict rule of \cref{sec:ren-network} discards the mode at $\tau_{\mathrm L}=1/b$; keeping it would give the singletons there). The region $C$ is neither. The interval $(1/\lambda_2(L_C),\lvert C\rvert/\cut(C))$ of \cref{conj:ren-lrg} is $(1/(b-1),1)$, which lies in $\tau_{\mathrm L}\ge1/b$. Here $N=b$, and for every constant $\mathsf c_0$ there is $b$ with $b-1>\mathsf c_0\ln b$, so no constant makes the threshold conjecture true. The whole graph has zero cut, so adding the missing vertex to $C$ raises $r_\star$ without bound: $C$ is not a local maximizer.
\end{proof}

The same failure for a level-one block of the nested medium with one site removed is a computation at one parameter value and is reported as such in \cref{chk:obj-lrg}.

\begin{proposition}[Mean specific heat of the infinite hierarchy]\label{prop:obj-specific-heat}
Let $\mathfrak q<1$, $\mathsf g_1^\infty=A_{\mathrm h}(1/(1-\mathfrak q)+1/(b-1))$ and $\mathsf z(\tau)=\sum_{l\ge1}(b-1)b^{-l}e^{-\tau\mathsf g_1^\infty\mathfrak q^{l-1}}$. As $n_{\mathrm H}\to\infty$, $N^{-1}\tr e^{-\tau L}$ converges to $\mathsf z(\tau)$ locally uniformly on $(0,\infty)$ together with its first two derivatives, so the entropy of $\varrho_\tau$ minus $\ln N$ converges to $\mathsf s(\tau)=\ln\mathsf z-\tau\mathsf z'/\mathsf z$ and its specific heat to $\mathsf C(\tau)=-\mathrm d\mathsf s/\mathrm d\ln\tau$. The limit satisfies
\[
\lim_{\tau\to\infty}\bigl[\mathsf s(\tau)-\mathsf s(\tau/\mathfrak q)\bigr]=\ln b,\qquad\text{so}\qquad\frac1{\ln(1/\mathfrak q)}\int_{\ln\tau}^{\ln\tau+\ln(1/\mathfrak q)}\mathsf C\mathrm d\ln\tau'\longrightarrow\frac{\ln b}{\ln(1/\mathfrak q)}=\frac1{z_{\mathrm H}}.
\]
\end{proposition}

\begin{proof}
By \cref{thm:obj-nested}(a) and the proof of its part (b), $N^{-1}\tr e^{-\tau L}=b^{-n_{\mathrm H}}+\sum_{l=1}^{n_{\mathrm H}}(b-1)b^{-l}e^{-\tau\mathsf g_l}$ with $\mathsf g_l=\mathsf g_1^\infty\mathfrak q^{l-1}-A_{\mathrm h}\mathfrak q^{n_{\mathrm H}}/(1-\mathfrak q)\le\mathsf g_1^\infty$. The $k$th derivative of the $l$th term is bounded by $(b-1)b^{-l}(\mathsf g_1^\infty)^k$ uniformly in $\tau$ and $n_{\mathrm H}$, and each term converges locally uniformly with its derivatives, so the series converges as stated. With $Z=\tr e^{-\tau L}$, $-\tr\varrho_\tau\ln\varrho_\tau=\tau\tr(L\varrho_\tau)+\ln Z=\ln Z-\tau Z'/Z$, which gives the entropy; $\mathsf z>0$ gives the convergence of the entropy and of the specific heat, which involves $\mathsf z''$.

Separating the first term of $\mathsf z(\tau/\mathfrak q)$ gives $\mathsf z(\tau/\mathfrak q)=\mathsf z(\tau)/b+\mathsf R(\tau)$ with $\mathsf R(\tau)=(b-1)b^{-1}e^{-\tau\mathsf g_1^\infty/\mathfrak q}$. Since $(\tau/\mathfrak q)\mathsf z'(\tau/\mathfrak q)=\tau\frac{\mathrm d}{\mathrm d\tau}\mathsf z(\tau/\mathfrak q)$,
\[
\mathsf s(\tau/\mathfrak q)-\mathsf s(\tau)=-\ln b+\ln\Bigl(1+\frac{b\mathsf R}{\mathsf z}\Bigr)-\frac{\tau b(\mathsf R'\mathsf z-\mathsf z'\mathsf R)}{(\mathsf z+b\mathsf R)\mathsf z}.
\]
Three bounds finish the proof. First, $\tau\lvert\mathsf R'\rvert=(\tau\mathsf g_1^\infty/\mathfrak q)\mathsf R$. Second,
\[
\tau\lvert\mathsf z'\rvert=\sum_l(b-1)b^{-l}\bigl(\tau\mathsf g_1^\infty\mathfrak q^{l-1}\bigr)e^{-\tau\mathsf g_1^\infty\mathfrak q^{l-1}}\le e^{-1},
\]
because $xe^{-x}\le e^{-1}$ and the weights sum to one. Third, for $\tau\mathsf g_1^\infty\ge1$ let $l$ be the least integer with $\tau\mathsf g_1^\infty\mathfrak q^{l-1}\le1$; then $l-1\le\ln(\tau\mathsf g_1^\infty)/\ln(1/\mathfrak q)+1$, and keeping only that term,
\[
\mathsf z(\tau)\ge(b-1)b^{-l}e^{-1}\ge(b-1)b^{-2}e^{-1}(\tau\mathsf g_1^\infty)^{-1/z_{\mathrm H}}.
\]
Hence $1/\mathsf z$ grows at most polynomially while $\mathsf R$ and $\tau\lvert\mathsf R'\rvert$ decay exponentially, so $b\mathsf R/\mathsf z\to0$ and
\[
\Bigl\lvert\frac{\tau b(\mathsf R'\mathsf z-\mathsf z'\mathsf R)}{(\mathsf z+b\mathsf R)\mathsf z}\Bigr\rvert\le\frac{b\tau\lvert\mathsf R'\rvert}{\mathsf z}+\frac{b\mathsf R}{e\mathsf z^2}\longrightarrow0.
\]
The integral of $\mathsf C$ over one period of $\ln\tau$ is $\mathsf s(\tau)-\mathsf s(\tau/\mathfrak q)$.
\end{proof}

The proposition fixes the mean of the specific heat over a period and says nothing about its variation within one. That the variation does not die out is observed in \cref{chk:obj-specific-heat} and is not proved.

\begin{proof}[Proof of \cref{prop:obj-clocks}]
The identities \eqref{eq:obj-nested-spectrum} hold at each time, so $\Pi(t)=\sum_{m=1}^{n_{\mathrm H}+1}(\varepsilon^\Pi(t)+\mathsf g_m(t))\mathsf E_m$ with $\mathsf g_m(t)$ as in \eqref{eq:obj-clock}: all $\Pi(t)$ lie in the commutative algebra spanned by the fixed projections $\mathsf E_m$, and the equation has a unique strong solution for locally integrable coefficients.

(a) Let $\Phi$ hold an orthonormal basis of $\operatorname{range}\mathsf P_l$, so $\Phi\Phi^\top=\mathsf P_l$. Since $\Pi(t)$ commutes with $\mathsf P_l$, $\Pi(t)\Phi=\mathsf P_l\Pi(t)\Phi=\Phi\Pi_\Phi(t)$, hence $\Phi^\top\Pi(t)=\Pi_\Phi(t)\Phi^\top$ and $\mathrm dy=-\Pi_\Phi(t)y\mathrm dt+\sqrt2\mathrm d\mathsf B'_t$ with $\mathsf B'=\Phi^\top\mathsf B$ a standard Brownian motion. This is a closed equation for $y$ with deterministic coefficients, whatever the law of $x_0$.

(b) Applying $\mathsf E_l$ to the equation, $\mathrm dx^{(l)}=-(\varepsilon^\Pi(t)+\mathsf g_l(t))x^{(l)}\mathrm dt+\sqrt2\mathsf E_l\mathrm d\mathsf B_t$, so
\[
x^{(l)}_t=e^{-\Theta_l(s,t)}x^{(l)}_s+\sqrt2\int_s^te^{-\Theta_l(r,t)}\mathsf E_l\mathrm d\mathsf B_r,
\]
where the integral is a centered Gaussian vector with covariance $2\int_s^te^{-2\Theta_l(r,t)}\mathrm dr\mathsf E_l$, independent of $(x_r)_{r\le s}$. This is the claim; when $x_s$ has no first moment the conditional expectation is understood in this sense. Each clock is absolutely continuous and additive, $\Theta_l(s,r)+\Theta_l(r,t)=\Theta_l(s,t)$, and strictly increasing in $t$ because its speed $\varepsilon^\Pi+\mathsf g_l$ is positive. If the leak were allowed to vanish on an interval, the clock of the population mean would stall there, and that of level $l$ wherever in addition the couplings of levels $l$ and above vanish.

(c) By (b), $x_t=\mathsf Ux_s+\zeta$ with $\mathsf U=\sum_le^{-\Theta_l(s,t)}\mathsf E_l$ and $\zeta$ Gaussian, independent of $x_s$, with covariance $\sum_l2\int_s^te^{-2\Theta_l(r,t)}\mathrm dr\mathsf E_l$. If $x_s$ is Gaussian with covariance $S=\sum_lv_l(s)\mathsf E_l$, the pair is jointly Gaussian, $\Cov(x_t)=\sum_lv_l(t)\mathsf E_l$ with $v_l(t)$ as stated, and
\[
\Cov(x_s\mid x_t)=S-S\mathsf U\Cov(x_t)^{-1}\mathsf US=S^{1/2}(I-\mathsf H)S^{1/2},\qquad\mathsf H=\sum_l\frac{e^{-2\Theta_l(s,t)}v_l(s)}{v_l(t)}\mathsf E_l=\sum_le^{-2\tilde\Theta_l}\mathsf E_l.
\]
The canonical correlations are the square roots of the eigenvalues of $\mathsf H$, that is $e^{-\tilde\Theta_l}$ with multiplicity $\rank\mathsf E_l$, and $\tilde\Theta_l>0$ because $v_l(t)>e^{-2\Theta_l(s,t)}v_l(s)$. Additivity of the clock gives
\[
e^{2\tilde\Theta_l}=1+\frac2{v_l(s)}\int_s^te^{2\Theta_l(s,r)}\mathrm dr.
\]
If $a^{\mathrm H}_l>a^{\mathrm H}_{l+1}$ almost everywhere on $[s,t]$, then $\mathsf g_l-\mathsf g_{l+1}=b^l(a^{\mathrm H}_l-a^{\mathrm H}_{l+1})>0$ there by \eqref{eq:obj-nested-spectrum}, so $\Theta_l(s,r)>\Theta_{l+1}(s,r)$ for $s<r\le t$, and with $v_l(s)\le v_{l+1}(s)$ the right-hand side is strictly larger for $l$ than for $l+1$: $\tilde\Theta_1>\dots>\tilde\Theta_{n_{\mathrm H}+1}$. Apply \cref{lem:obj-compression} to $(x_s,x_t)$. For $M=b^{n_{\mathrm H}-l}$ the ordered exponents satisfy $\theta_M=\tilde\Theta_{l+1}<\theta_{M+1}=\tilde\Theta_l$, the span of the first $M$ eigenvectors of $\mathsf H$ is $\bigoplus_{m>l}\operatorname{range}\mathsf E_m=\operatorname{range}\mathsf P_l$, and $S^{1/2}$ maps this space onto itself. The argument of \cref{thm:obj-selection}(b) then shows that the frames with range $\operatorname{range}\mathsf P_l$ are the only minimizers exactly when $\tilde\Theta_{l+1}<\varsigma(c)<\tilde\Theta_l$. The hypothesis $v_l(s)\le v_{l+1}(s)$ holds for an isotropic law and for the stationary law of any nested medium with decreasing level couplings, where $v_l=1/(\varepsilon^\Pi+\mathsf g_l)$; in a constant medium started there, $\tilde\Theta_l=(\varepsilon^\Pi+\mathsf g_l)(t-s)$ and the statement is \cref{thm:obj-nested}(a).
\end{proof}

\begin{remark}[What blocks a general non-stationary statement]\label{rem:obj-nonstationary}
\Cref{lem:obj-compression} holds for any jointly Gaussian pair, so for any linear Gaussian dynamics and any two times $s<t$ the frames that minimize the criterion are those whose image under $\Cov(x_s)^{1/2}$ spans the leading eigenvectors of $\mathsf H(s,t)$, and the selected rank is the number of canonical correlations above $e^{-\varsigma(c)}$. What \cref{prop:obj-clocks} adds on a fixed tree, and what fails in general, is that these frames are the same for all pairs of times and are closed. Two exact computations locate the failure. For a fixed frame $\Phi$ with complement $\Phi_\perp$ and continuous symmetric $\Pi(t)$, eliminating $\Phi_\perp^\top x$ gives
\[
\dot y(t)=-\Pi_\Phi(t)y(t)+\int_s^t\Gamma(t,r)y(r)\mathrm dr+\mathfrak f(t),\qquad\Gamma(t,r)=\Pi_{\perp\Phi}(t)^\top\mathsf U_\perp(t,r)\Pi_{\perp\Phi}(r),
\]
with $\mathsf U_\perp$ the propagator of $\dot w=-\Pi_{\perp\perp}(t)w$ and $\mathfrak f$ a forcing that does not depend on $y$; since $\Gamma(t,t)=\Pi_{\perp\Phi}(t)^\top\Pi_{\perp\Phi}(t)$, the kernel vanishes identically if and only if $\operatorname{range}\Phi$ is invariant under every $\Pi(t)$. For a differentiable frame $\Phi(t)$ whose range is invariant under $\Pi(t)$ at each time, $y=\Phi(t)^\top x$ obeys $\mathrm dy=(\dot\Phi^\top\Phi-\Pi_\Phi)y\mathrm dt+\dot\Phi^\top\Phi_\perp\Phi_\perp^\top x\mathrm dt+\sqrt2\Phi^\top\mathrm d\mathsf B_t$, and $\dot\Phi^\top\Phi_\perp=0$ for all $t$ if and only if the projection $\Phi\Phi^\top$ is constant. Exact closure at all times therefore requires a common invariant frame, which is the situation of \cref{prop:obj-clocks}. For a nonreversible stationary drift $K$ with covariance $\Sigma$, $\mathsf H(\tau)=e^{-\tilde{\mathsf K}^\top\tau}e^{-\tilde{\mathsf K}\tau}$ with $\tilde{\mathsf K}=\Sigma^{-1/2}K\Sigma^{1/2}$; its eigenvectors are those of $\tilde{\mathsf K}+\tilde{\mathsf K}^\top$ at every horizon when $\tilde{\mathsf K}$ is normal, and we have no horizon-independent description of them otherwise. The computations are exact. That they are the only obstruction is not claimed, and a scale-time law for these cases is part of \cref{open:obj-scale}.
\end{remark}
